\section{Background}
\label{sec:background}

% THIS IS THE SECTION WHERE WE GIVE AWAY THE TRICK.
% The nice part of the writing in PLDI is that by the 
% end of the background section, you sort of already 'get'
% the paper. The theory is really just there to confirm
% the intuition.

% TRICK: CERTAIN KINDS OF PROGRAMS CAN COMPILE TO 
% K-LINEAR MAPS.

The most important technical idea in this paper is that $\cj$\! 
programs compile to $\rd{k}$-linear maps. Because these $\rd{k}$-linear maps
are differentiable, we can embed them within a transformer.

\paragraph{Linear Algebra} Linear maps model additive interactions between neurons. 
However, $\rd{k}$-linear maps additionally model multiplicative interactions, while ensuring
linearity in each argument.
Consider the difference below, between the linear map $\rd{f}$ and the 
$\rd{2}$-linear map $\rd{g}$.

\[\rd{f}\Bigl(\rd{\begin{bmatrix}x_1\\x_2\end{bmatrix}}\Bigr) = 
    \rd{\begin{bmatrix}2x_1+3x_2\\4x_1\end{bmatrix}}\quad\quad 
\rd{g}\Bigl(\rd{\begin{bmatrix}x_1\\x_2\end{bmatrix}}, \rd{\begin{bmatrix}y_1\\y_2\end{bmatrix}}\Bigr)=
    \rd{\begin{bmatrix}2x_1y_2+3x_1y_1\\4x_2y_1\end{bmatrix}}
\]
\vspace{.2em}

Multiplicative interactions (e.g., $\rd{2x_1y_2}$) can implement more intricate behaviors,
like gating. For example, the output of $\rd{g}$ is gated by its first input $\bv{x_1}{x_2}$. When $\rd{x_1}$
is $\rd{0}$, it turns off the activity of the first output neuron. Likewise when $\rd{x_2}$
is $\rd{0}$, it turns off the activity of the second output neuron. This style of gating
is how we will implement conditionals using $\rd{k}$-linear maps, including for the conditional
reverse program shown in Fig. \ref{fig:conditional_reverse}.

\paragraph{Linear Types}
% Reader must know what a typing judgment is here.
Most programming languages do not compile to $\rd{k}$-linear maps. However, some languages
with linear types can. These types annotate programs with information about their behavior,
including constraints the behavior satisfies.\footnote{In this paper, a program's behavior
is the value produced by evaluating it.} Consider the difference below, between the type annotation for the conditional on the 
left and right.

\vspace{-.6em}
\[
{\text{\textcolor{ForestGreen}{\Large \cmark}}}\;\;\bind{\nb{x}}{\nb{\t{Bool}}},\bind{\nb{y}}{\nb{\t{Bool}}} \- 
    \t{if}\,\nb{x}\,\t{then}\,\nb{y}\,\t{else}\,\nb{y} : \nb{\t{Bool}}
\quad\!\!{\large \boldsymbol{\mid}}\!\!\quad
\bind{\nb{x}}{\nb{\t{Bool}}},\bind{\nb{y}}{\nb{\t{Bool}}} \- 
    \t{if}\,\nb{x}\,\t{then}\,\nb{x}\,\t{else}\,\nb{x} : \nb{\t{Bool}}
    \;\;
    {\text{\textcolor{Red}{\Large \xmark}}}
\]

Both are legitimate type annotations. They annotate the free variables $\nb{x}$ and $\nb{y}$ 
with types that make sense (e.g., $\nb{x}$ must be $\nb{\t{Bool}}$ for the conditional to have
behavior), and also the type of value they evaluate to---when free variables are instantiated by
values of their annotated type. But only the type annotation on the left is a linear type 
annotation: it ensures that evaluation of a program uses variables \textit{exactly once}.
It uses $\nb{x}$ when checking the condition and $\nb{y}$ when either branch executes.
For the type annotation on the right, it uses $\nb{x}$ both in checking the condition and
when either branch executes $\nb{x}$, and does not use $\nb{y}$ at all.

A remarkable fact of linear types is that these syntactic constraints on permissible programs
correspond to the constraints imposed on $\rd{k}$-linear maps. This has been known for decades
in programming language semantics \citep{bierman1995categorical,mellies2009categorical}, but only recently used to design programming languages
whose programs compile to neural networks \citep{velez2026compiling,velez2026compiling2}.
For example, the \textit{linear} conditional above compiles to the following $\rd{2}$-linear map, which
can exactly implement the behavior of the program.\footnote{Assuming a typical one-hot encoding of booleans.}

\[
\llbracket\bind{\nb{x}}{\nb{\t{Bool}}},\bind{\nb{y}}{\nb{\t{Bool}}} \- 
    \t{if}\,\nb{x}\,\t{then}\,\nb{y}\,\t{else}\,\nb{y} : \nb{\t{Bool}}\rrbracket
    \Bigl(\rd{\begin{bmatrix}x_1 \\ x_2\end{bmatrix}}, 
    \rd{\begin{bmatrix}y_1 \\ y_2\end{bmatrix}}\Bigr) = 
    \rd{\begin{bmatrix}x_1y_1+x_2y_1 \\ x_1y_2+x_2y_2\end{bmatrix}}
\]

The goal of the following technical material is to formally establish this correspondence across all
programs in $\cj$\!, and to develop a compiler which computes the correspondence for us.
However, the most important intuitions are precisely those covered here in the background.
